\documentclass[aspectratio=169,11pt]{beamer}
% Week 2 Practice — Python warm-ups for data
% Course: AI Frameworks (ages 16--18)

% Shared Beamer preamble for AI Frameworks
% Audience: college students ages 16--18 (intuition over math)
% Include with: % Shared Beamer preamble for AI Frameworks
% Audience: college students ages 16--18 (intuition over math)
% Include with: % Shared Beamer preamble for AI Frameworks
% Audience: college students ages 16--18 (intuition over math)
% Include with: \input{preamble}

\usetheme{Madrid}
\usecolortheme{default}

% Clean palette: deep slate + teal accent (distinct from Java navy decks)
\definecolor{LectureNavy}{HTML}{1A365D}
\definecolor{LectureAccent}{HTML}{319795}
\definecolor{LectureGray}{HTML}{4A5568}
\definecolor{CodeBg}{HTML}{F7FAFC}
\definecolor{AlertBg}{HTML}{FFF5F5}
\definecolor{TryBg}{HTML}{F0FFF4}
\definecolor{QuizBg}{HTML}{EBF8FF}
\definecolor{AnswerKeyBg}{HTML}{FFFFF0}
\definecolor{AnswerGold}{HTML}{B7791F}
\definecolor{KeywordBlue}{HTML}{2B6CB0}
\definecolor{StringGreen}{HTML}{276749}
\definecolor{CommentGray}{HTML}{718096}

\setbeamercolor{structure}{fg=LectureNavy}
\setbeamercolor{palette primary}{bg=LectureNavy,fg=white}
\setbeamercolor{palette secondary}{bg=LectureAccent,fg=white}
\setbeamercolor{palette tertiary}{bg=LectureGray,fg=white}
% Madrid title box is dark --- title/subtitle must be light for contrast
\setbeamercolor{title}{fg=white,bg=LectureNavy}
\setbeamercolor{subtitle}{fg=white,bg=LectureNavy}
\setbeamercolor{author}{fg=LectureNavy}
\setbeamercolor{date}{fg=LectureGray}
\setbeamercolor{institute}{fg=LectureGray}
\setbeamercolor{frametitle}{bg=LectureNavy,fg=white}
\setbeamercolor{block title}{bg=LectureNavy,fg=white}
\setbeamercolor{block body}{bg=CodeBg,fg=black}
\setbeamercolor{block title alerted}{bg=LectureAccent,fg=white}
\setbeamercolor{block body alerted}{bg=TryBg,fg=black}
\setbeamercolor{block title example}{bg=LectureGray,fg=white}
\setbeamercolor{block body example}{bg=AlertBg,fg=black}

\setbeamertemplate{navigation symbols}{}
\setbeamertemplate{itemize items}[circle]
\setbeamertemplate{enumerate items}[default]

\usepackage[T1]{fontenc}
\usepackage[utf8]{inputenc}
\usepackage{lmodern}
\usepackage{booktabs}
\usepackage{tabularx}
\usepackage{graphicx}
\usepackage{hyperref}
\usepackage{listings}
\usepackage{xcolor}
\usepackage{tikz}
\usetikzlibrary{positioning,arrows.meta,shapes.geometric}

% listings: Python without shell-escape
\lstdefinestyle{pythonstyle}{
    language=Python,
    basicstyle=\ttfamily\small,
    keywordstyle=\color{KeywordBlue}\bfseries,
    commentstyle=\color{CommentGray}\itshape,
    stringstyle=\color{StringGreen},
    numbers=left,
    numberstyle=\tiny\color{CommentGray},
    numbersep=6pt,
    frame=single,
    rulecolor=\color{LectureGray!40},
    backgroundcolor=\color{CodeBg},
    breaklines=true,
    showstringspaces=false,
    tabsize=4,
    captionpos=b,
}
\lstset{style=pythonstyle}

\newcommand{\pycode}[1]{{\ttfamily\small\detokenize{#1}}}
\newcommand{\trythis}[1]{%
    \begin{alertblock}{Try this}
    #1
    \end{alertblock}%
}
\newcommand{\commonmistake}[1]{%
    \begin{exampleblock}{Common mistake}
    #1
    \end{exampleblock}%
}
\newcommand{\remember}[1]{%
    \begin{block}{Remember}
    #1
    \end{block}%
}

% Formative in-class checks: question slide, then a dedicated Answer slide
\newcommand{\quizpause}{%
    \par\vspace{0.5em}
    {\small\textit{Pause --- answer (clicker / raise hand / neighbour) before the next slide.}}%
}
\newcommand{\quizbox}[1]{%
    \setbeamercolor{block title}{bg=KeywordBlue,fg=white}%
    \setbeamercolor{block body}{bg=QuizBg,fg=black}%
    \begin{block}{Quiz}
    #1
    \end{block}%
    \setbeamercolor{block title}{bg=LectureNavy,fg=white}%
    \setbeamercolor{block body}{bg=CodeBg,fg=black}%
}
\newcommand{\answerbox}[1]{%
    \setbeamercolor{block title}{bg=AnswerGold,fg=white}%
    \setbeamercolor{block body}{bg=AnswerKeyBg,fg=black}%
    \begin{block}{Answer}
    #1
    \end{block}%
    \setbeamercolor{block title}{bg=LectureNavy,fg=white}%
    \setbeamercolor{block body}{bg=CodeBg,fg=black}%
}

% Empty author/date placeholders for the lecturer to fill
\author{}
\date{}


\usetheme{Madrid}
\usecolortheme{default}

% Clean palette: deep slate + teal accent (distinct from Java navy decks)
\definecolor{LectureNavy}{HTML}{1A365D}
\definecolor{LectureAccent}{HTML}{319795}
\definecolor{LectureGray}{HTML}{4A5568}
\definecolor{CodeBg}{HTML}{F7FAFC}
\definecolor{AlertBg}{HTML}{FFF5F5}
\definecolor{TryBg}{HTML}{F0FFF4}
\definecolor{QuizBg}{HTML}{EBF8FF}
\definecolor{AnswerKeyBg}{HTML}{FFFFF0}
\definecolor{AnswerGold}{HTML}{B7791F}
\definecolor{KeywordBlue}{HTML}{2B6CB0}
\definecolor{StringGreen}{HTML}{276749}
\definecolor{CommentGray}{HTML}{718096}

\setbeamercolor{structure}{fg=LectureNavy}
\setbeamercolor{palette primary}{bg=LectureNavy,fg=white}
\setbeamercolor{palette secondary}{bg=LectureAccent,fg=white}
\setbeamercolor{palette tertiary}{bg=LectureGray,fg=white}
% Madrid title box is dark --- title/subtitle must be light for contrast
\setbeamercolor{title}{fg=white,bg=LectureNavy}
\setbeamercolor{subtitle}{fg=white,bg=LectureNavy}
\setbeamercolor{author}{fg=LectureNavy}
\setbeamercolor{date}{fg=LectureGray}
\setbeamercolor{institute}{fg=LectureGray}
\setbeamercolor{frametitle}{bg=LectureNavy,fg=white}
\setbeamercolor{block title}{bg=LectureNavy,fg=white}
\setbeamercolor{block body}{bg=CodeBg,fg=black}
\setbeamercolor{block title alerted}{bg=LectureAccent,fg=white}
\setbeamercolor{block body alerted}{bg=TryBg,fg=black}
\setbeamercolor{block title example}{bg=LectureGray,fg=white}
\setbeamercolor{block body example}{bg=AlertBg,fg=black}

\setbeamertemplate{navigation symbols}{}
\setbeamertemplate{itemize items}[circle]
\setbeamertemplate{enumerate items}[default]

\usepackage[T1]{fontenc}
\usepackage[utf8]{inputenc}
\usepackage{lmodern}
\usepackage{booktabs}
\usepackage{tabularx}
\usepackage{graphicx}
\usepackage{hyperref}
\usepackage{listings}
\usepackage{xcolor}
\usepackage{tikz}
\usetikzlibrary{positioning,arrows.meta,shapes.geometric}

% listings: Python without shell-escape
\lstdefinestyle{pythonstyle}{
    language=Python,
    basicstyle=\ttfamily\small,
    keywordstyle=\color{KeywordBlue}\bfseries,
    commentstyle=\color{CommentGray}\itshape,
    stringstyle=\color{StringGreen},
    numbers=left,
    numberstyle=\tiny\color{CommentGray},
    numbersep=6pt,
    frame=single,
    rulecolor=\color{LectureGray!40},
    backgroundcolor=\color{CodeBg},
    breaklines=true,
    showstringspaces=false,
    tabsize=4,
    captionpos=b,
}
\lstset{style=pythonstyle}

\newcommand{\pycode}[1]{{\ttfamily\small\detokenize{#1}}}
\newcommand{\trythis}[1]{%
    \begin{alertblock}{Try this}
    #1
    \end{alertblock}%
}
\newcommand{\commonmistake}[1]{%
    \begin{exampleblock}{Common mistake}
    #1
    \end{exampleblock}%
}
\newcommand{\remember}[1]{%
    \begin{block}{Remember}
    #1
    \end{block}%
}

% Formative in-class checks: question slide, then a dedicated Answer slide
\newcommand{\quizpause}{%
    \par\vspace{0.5em}
    {\small\textit{Pause --- answer (clicker / raise hand / neighbour) before the next slide.}}%
}
\newcommand{\quizbox}[1]{%
    \setbeamercolor{block title}{bg=KeywordBlue,fg=white}%
    \setbeamercolor{block body}{bg=QuizBg,fg=black}%
    \begin{block}{Quiz}
    #1
    \end{block}%
    \setbeamercolor{block title}{bg=LectureNavy,fg=white}%
    \setbeamercolor{block body}{bg=CodeBg,fg=black}%
}
\newcommand{\answerbox}[1]{%
    \setbeamercolor{block title}{bg=AnswerGold,fg=white}%
    \setbeamercolor{block body}{bg=AnswerKeyBg,fg=black}%
    \begin{block}{Answer}
    #1
    \end{block}%
    \setbeamercolor{block title}{bg=LectureNavy,fg=white}%
    \setbeamercolor{block body}{bg=CodeBg,fg=black}%
}

% Empty author/date placeholders for the lecturer to fill
\author{}
\date{}


\usetheme{Madrid}
\usecolortheme{default}

% Clean palette: deep slate + teal accent (distinct from Java navy decks)
\definecolor{LectureNavy}{HTML}{1A365D}
\definecolor{LectureAccent}{HTML}{319795}
\definecolor{LectureGray}{HTML}{4A5568}
\definecolor{CodeBg}{HTML}{F7FAFC}
\definecolor{AlertBg}{HTML}{FFF5F5}
\definecolor{TryBg}{HTML}{F0FFF4}
\definecolor{QuizBg}{HTML}{EBF8FF}
\definecolor{AnswerKeyBg}{HTML}{FFFFF0}
\definecolor{AnswerGold}{HTML}{B7791F}
\definecolor{KeywordBlue}{HTML}{2B6CB0}
\definecolor{StringGreen}{HTML}{276749}
\definecolor{CommentGray}{HTML}{718096}

\setbeamercolor{structure}{fg=LectureNavy}
\setbeamercolor{palette primary}{bg=LectureNavy,fg=white}
\setbeamercolor{palette secondary}{bg=LectureAccent,fg=white}
\setbeamercolor{palette tertiary}{bg=LectureGray,fg=white}
% Madrid title box is dark --- title/subtitle must be light for contrast
\setbeamercolor{title}{fg=white,bg=LectureNavy}
\setbeamercolor{subtitle}{fg=white,bg=LectureNavy}
\setbeamercolor{author}{fg=LectureNavy}
\setbeamercolor{date}{fg=LectureGray}
\setbeamercolor{institute}{fg=LectureGray}
\setbeamercolor{frametitle}{bg=LectureNavy,fg=white}
\setbeamercolor{block title}{bg=LectureNavy,fg=white}
\setbeamercolor{block body}{bg=CodeBg,fg=black}
\setbeamercolor{block title alerted}{bg=LectureAccent,fg=white}
\setbeamercolor{block body alerted}{bg=TryBg,fg=black}
\setbeamercolor{block title example}{bg=LectureGray,fg=white}
\setbeamercolor{block body example}{bg=AlertBg,fg=black}

\setbeamertemplate{navigation symbols}{}
\setbeamertemplate{itemize items}[circle]
\setbeamertemplate{enumerate items}[default]

\usepackage[T1]{fontenc}
\usepackage[utf8]{inputenc}
\usepackage{lmodern}
\usepackage{booktabs}
\usepackage{tabularx}
\usepackage{graphicx}
\usepackage{hyperref}
\usepackage{listings}
\usepackage{xcolor}
\usepackage{tikz}
\usetikzlibrary{positioning,arrows.meta,shapes.geometric}

% listings: Python without shell-escape
\lstdefinestyle{pythonstyle}{
    language=Python,
    basicstyle=\ttfamily\small,
    keywordstyle=\color{KeywordBlue}\bfseries,
    commentstyle=\color{CommentGray}\itshape,
    stringstyle=\color{StringGreen},
    numbers=left,
    numberstyle=\tiny\color{CommentGray},
    numbersep=6pt,
    frame=single,
    rulecolor=\color{LectureGray!40},
    backgroundcolor=\color{CodeBg},
    breaklines=true,
    showstringspaces=false,
    tabsize=4,
    captionpos=b,
}
\lstset{style=pythonstyle}

\newcommand{\pycode}[1]{{\ttfamily\small\detokenize{#1}}}
\newcommand{\trythis}[1]{%
    \begin{alertblock}{Try this}
    #1
    \end{alertblock}%
}
\newcommand{\commonmistake}[1]{%
    \begin{exampleblock}{Common mistake}
    #1
    \end{exampleblock}%
}
\newcommand{\remember}[1]{%
    \begin{block}{Remember}
    #1
    \end{block}%
}

% Formative in-class checks: question slide, then a dedicated Answer slide
\newcommand{\quizpause}{%
    \par\vspace{0.5em}
    {\small\textit{Pause --- answer (clicker / raise hand / neighbour) before the next slide.}}%
}
\newcommand{\quizbox}[1]{%
    \setbeamercolor{block title}{bg=KeywordBlue,fg=white}%
    \setbeamercolor{block body}{bg=QuizBg,fg=black}%
    \begin{block}{Quiz}
    #1
    \end{block}%
    \setbeamercolor{block title}{bg=LectureNavy,fg=white}%
    \setbeamercolor{block body}{bg=CodeBg,fg=black}%
}
\newcommand{\answerbox}[1]{%
    \setbeamercolor{block title}{bg=AnswerGold,fg=white}%
    \setbeamercolor{block body}{bg=AnswerKeyBg,fg=black}%
    \begin{block}{Answer}
    #1
    \end{block}%
    \setbeamercolor{block title}{bg=LectureNavy,fg=white}%
    \setbeamercolor{block body}{bg=CodeBg,fg=black}%
}

% Empty author/date placeholders for the lecturer to fill
\author{}
\date{}


\title{Week 2 Practice: Python Warm-ups}
\subtitle{AI Frameworks}

\begin{document}

    \begin{frame}
        \titlepage
    \end{frame}

    \begin{frame}{Practice goals (80 minutes)}
        By the end of this session you should have:
        \begin{enumerate}
            \item Completed short exercises on lists and strings
            \item Written functions for \textbf{mean}, \textbf{min}, and \textbf{max} (no pandas)
            \item (Optional) Read a tiny CSV with the \pycode{csv} module
            \item A clear note: ``next week pandas will make tables easier''
        \end{enumerate}
    \end{frame}

    \begin{frame}{Agenda}
        \begin{enumerate}
            \item Warm-up drills (20--25 min)
            \item Mean / min / max functions (25--30 min)
            \item Optional: tiny CSV with \pycode{csv} (15 min)
            \item Quiz + done checklist (10--15 min)
        \end{enumerate}
    \end{frame}

%======================================================================
    \section{Warm-ups}
%======================================================================

    \begin{frame}[fragile]{Drill 1 --- lists of numbers}
        \begin{lstlisting}
nums = [4, 1, 9, 2, 6]

# TODO: print the length, the first item, the last item
# TODO: append 0 and print the new list
# TODO: create a new list with only values > 5 (use a loop)
        \end{lstlisting}
        \vspace{0.3em}
        Check results by hand before asking for help.
    \end{frame}

    \begin{frame}[fragile]{Drill 2 --- strings as sequences}
        \begin{lstlisting}
word = "Frameworks"

print(len(word))
print(word[0], word[-1])

# TODO: count how many letters are lowercase 'a' (use a loop)
# TODO: build a new string with only the first 4 characters (slice)
        \end{lstlisting}
    \end{frame}

    \begin{frame}[fragile]{Drill 3 --- a list of dicts (mini ``table'')}
        \begin{lstlisting}
rows = [
    {"name": "Ada", "score": 90},
    {"name": "Alan", "score": 80},
    {"name": "Grace", "score": 95},
]

# TODO: print each name
# TODO: compute the average score with a loop
        \end{lstlisting}
        \vspace{0.3em}
        This is the awkward version of a table---pandas will feel nicer next week.
    \end{frame}

%======================================================================
    \section{Mean, min, max}
%======================================================================

    \begin{frame}{Task --- write three functions}
        Without using Python's \pycode{min}, \pycode{max}, or statistics libraries:
        \begin{itemize}
            \item \pycode{my_mean(values)}
            \item \pycode{my_min(values)}
            \item \pycode{my_max(values)}
        \end{itemize}
        \vspace{0.5em}
        Assume \pycode{values} is a non-empty list of numbers.
    \end{frame}

    \begin{frame}[fragile]{Starter template}
        \begin{lstlisting}
def my_mean(values):
    total = 0
    for v in values:
        total = total + v
    return total / len(values)

def my_min(values):
    pass  # TODO

def my_max(values):
    pass  # TODO

data = [5, 2, 9, 2, 7]
print(my_mean(data), my_min(data), my_max(data))
        \end{lstlisting}
    \end{frame}

    \begin{frame}{Check yourself}
        \begin{itemize}
            \item For \pycode{[5, 2, 9, 2, 7]}: mean should be 5.0; min 2; max 9
            \item Try a second list to be sure
            \item Optional stretch: what should happen on an empty list? (discuss---don't crash silently)
        \end{itemize}
        \vspace{0.4em}
        \trythis{After yours work, compare one result to Python's built-in \pycode{min}/\pycode{max}.}
    \end{frame}

%======================================================================
    \section{Optional CSV}
%======================================================================

    \begin{frame}{Optional --- peek at CSV}
        Create a file \pycode{tiny.csv} (or use one provided):
        \vspace{0.3em}
        {\small
        \texttt{name,score}\\
        \texttt{Ada,90}\\
        \texttt{Alan,80}\\
        \texttt{Grace,95}
        }
    \end{frame}

    \begin{frame}[fragile]{Read with the csv module}
        \begin{lstlisting}
import csv

scores = []
with open("tiny.csv", "r", newline="") as f:
    reader = csv.DictReader(f)
    for row in reader:
        scores.append(float(row["score"]))

print(scores)
print(my_mean(scores))
        \end{lstlisting}
        \vspace{0.3em}
        Next week: \pycode{pd.read_csv("tiny.csv")} --- one clear line.
    \end{frame}

    \begin{frame}{If the optional part is slow}
        \begin{itemize}
            \item Skip CSV and finish mean/min/max solidly
            \item Or use Colab's file upload / a path the instructor shares
        \end{itemize}
        \vspace{0.4em}
        \commonmistake{Spending all remaining time fighting file paths.
        Core goal today = functions on lists.}
    \end{frame}

%--- Quiz ---
    \begin{frame}[fragile]{Quiz --- Practice check}
        \begin{lstlisting}
def my_max(values):
    best = values[0]
    for v in values:
        if v > best:
            best = v
    return best

print(my_max([3, 8, 2]))
        \end{lstlisting}
        \quizbox{What is printed?}
        \begin{enumerate}
            \item \pycode{3}
            \item \pycode{8}
            \item \pycode{2}
            \item \pycode{[3, 8, 2]}
        \end{enumerate}
        \quizpause
    \end{frame}

    \begin{frame}{Answer --- Practice check}
        \answerbox{\textbf{2.} \pycode{8} --- the loop keeps the largest value seen so far.}
    \end{frame}

%======================================================================
    \section{Done checklist}
%======================================================================

    \begin{frame}{Done checklist}
        \begin{enumerate}
            \item Warm-up drills on lists / strings attempted in a notebook
            \item \pycode{my\_mean}, \pycode{my\_min}, \pycode{my\_max} work on a sample list
            \item (Optional) CSV scores loaded and averaged
            \item You can explain list vs dict in one sentence each
        \end{enumerate}
    \end{frame}

    \begin{frame}{Before you leave}
        \begin{itemize}
            \item Save \pycode{week02.ipynb}
            \item Keep your functions---handy to compare with pandas later
        \end{itemize}
        \vspace{0.6em}
        \begin{block}{Next week}
            Lecture: pandas DataFrames\\
            Practice: explore a friendly CSV with \pycode{head}, filters, simple questions
        \end{block}
    \end{frame}

\end{document}
